\section{Field transformations}

Most concepts developed in this paper
are expected to be widely applicable, but as explained above,
the case of immediate interest is lattice QCD.
In the following, the gauge group is therefore taken to
be $\Group$.
Since the quarks will play a spectator r\^ole,
they will be added to the theory only in sect.~6, where the proposed
simulation algorithm for lattice QCD is discussed.

\subsection{Field space}

The lattice theory is set up on a finite
hypercubic lattice $\Lambda$ with periodic boundary conditions.
For notational convenience, the lattice spacing is set to unity.
As usual, the gauge field variables $U(x,\mu)\in\Group$
are assumed to
reside on the links $(x,\mu)$ of the lattice
(where $x\in\Lambda$ and $\mu=0,\ldots,3$).
The expectation value of any observable
$\mathcal{O}(U)$ is then given by the functional integral
\begin{equation}
  \langle\mathcal{O}\rangle=
  \frac{1}{\mathcal{Z}}\int\rmD[U]\,\mathcal{O}(U)\,\rme^{-S(U)}
\end{equation}
over the space of gauge fields. In this expression,
$S(U)$ denotes the gauge-field action,
$\mathcal{Z}$ the partition function and
$\rmD[U]$ the product of the normalized
$\Group$-invariant integration measures of the link variables
$U(x,\mu)$.

From a purely mathematical point of view, the space of
lattice gauge fields
is a power of $\Group$ and therefore a compact connected manifold.
Field transformations
are invertible maps of this manifold onto itself.
Such transformations will always be required to be
differentiable
in both directions and orientation-preserving
(here and below, ``differentiable'' means
``infinitely often continuously differentiable'').


\subsection{Right-invariant differential operators}

Since the link variables $U(x,\mu)$ take values in a Lie group,
it is natural to express differentiations with respect
to them through a basis
$\du^a_{x,\mu}$, $a=1,\ldots,8$, of
differential operators that are invariant under the right-action
of the group.
The action of these operators
on a differentiable function $f(U)$ of the gauge field
is given by
\begin{equation}
  \du_{x,\mu}^af(U)=\left.{\rmd\over\rmd t}
  f(U_t)\right|_{t=0},\qquad\quad
  U_t(y,\nu)=
  \begin{cases}
    \rme^{tT^a}U(x,\mu) & \text{if $(y,\nu)=(x,\mu)$,}\\[1.0ex]
    U(y,\nu)            & \text{otherwise,}
  \end{cases}
\end{equation}
where $T^a$ are the $\Group$ generators (see appendix~A).
In particular,
$\du^a_{x,\mu}$ transforms according to the adjoint
representation under the left-action of $\Group$.

The operators $\du^a_{x,\mu}$ go along with
a basis of 1-forms on the field manifold,
\begin{equation}
  \theta^a_{x,\mu}=-2\,\tr\{\rmd U(x,\mu)U(x,\mu)^{-1}T^a\},
\end{equation}
such that
\begin{equation}
  \rmd f(U)=\sum_{x,\mu}\theta^a_{x,\mu}\du^a_{x,\mu}f(U)
\end{equation}
for all functions $f(U)$.


\subsection{Jacobian matrix}

For any given gauge field $V(y,\nu)$,
the transformation (1.1) produces another field
$U(x,\mu)=[\trans(V)](x,\mu)$.
When considering such transformations, one
needs to distinguish differentiations with respect to $V$
from those with respect $U$. The associated 1-forms must
also be distinguished. In the following, all symbols
carrying a hat represent quantities and operations
referring to $V$.

The Jacobian matrix
\begin{equation}
  [\transs(V)](x,\mu;y,\nu)^{ab}=
  -2\,\tr\{\hat{\du}^b_{y,\nu}U(x,\mu)U(x,\mu)^{-1}T^a\}
\end{equation}
can be considered to be the kernel of a linear operator acting
on link fields with
values in $\Lie$. In particular,
\begin{equation}
  \theta^a_{x,\mu}=\sum_{y,\nu}[\transs(V)](x,\mu;y,\nu)^{ab}
  \hat{\theta}^b_{y,\nu}.
\end{equation}
Since the functional integration measure $\rmD[U]$
is proportional to the maximal product of these
1-forms, it follows that
\begin{equation}
  \rmD[U]=\rmD[V]\det\transs(V).
\end{equation}
The Jacobian of the map (1.1) is thus $\det\transs(V)$.

If the transformation satisfies
\begin{equation}
  S(\trans(V))-\ln\det\transs(V)=\text{constant},
\end{equation}
the substitution $U\to V$ of the integration variables
in the functional integral
maps the theory to the trivial one where the link variables
are completely decoupled from one another.
The expectation values (2.1) are then given by
\begin{equation}
  \langle\mathcal{O}\rangle=
  \int\rmD[V]\,\mathcal{O}(\trans(V)).
\end{equation}
Such trivializing maps thus contain the entire
dynamics of the theory.

Although the remark is likely to remain an academic one,
an intriguing observation is that the integral (2.9)
can be simulated simply by generating
uniformly distributed random gauge fields. Subsequent
field configurations are uncorrelated in this case
and there are therefore no autocorrelations in
the data series for the observables $\mathcal{O}$.


\subsection{Transformation behaviour of the HMC algorithm}

The HMC algorithm~\cite{HMC} operates on the phase space
associated to the field manifold.
In particular, the transition $V\to V'$ in fig.~1 requires
the equations of motion derived from the Hamilton function
\begin{equation}
  \hat{H}(\hat{\mom},V)=
  \frac{1}{2}(\hat{\mom},\hat{\mom})+
  S(\trans(V))-\ln\det\transs(V)
\end{equation}
to be integrated,
where $\hat{\mom}(x,\mu)\in\Lie$ is the canonical momentum of $V(x,\mu)$.

Although the complete update algorithm for the field $U$
described by fig.~1 looks different,
it is in fact equivalent to the HMC algorithm
with a non-standard Hamilton function.
The equivalence can be established by noting
that the transformation from $V$ to $U$
preserves the symplectic 2-form
\begin{equation}
  \hat{\Omega}=
  \sum_{x,\mu}\rmd\{\hat{\mom}^a(x,\mu)\hat{\theta}^a_{x,\mu}\},
\end{equation}
i.e.\ $\Omega=\hat{\Omega}$,
if the momenta of the fields are transformed according to
\begin{equation}
  \hat{\mom}^b(y,\nu)
  =\sum_{x,\mu}[\transs(V)](x,\mu;y,\nu)^{ab}\mom^a(x,\mu).
\end{equation}
The evolution of the transformed fields
is then governed by the Hamilton function
\begin{equation}
  H(\mom,U)=
  \frac{1}{2}(\mom,K(U)\mom)+
  S(U)-\ln\det\transs(V),
\end{equation}
where
\begin{equation}
  K(U)=\transs(V)^T\transs(V),\qquad V=\trans^{-1}(U).
\end{equation}
Note that the Jacobian in eq.~(2.13)
cancels when the momenta are integrated over in the
functional integral.
The algorithm outlined in fig.~1 thus amounts
to applying the HMC algorithm with
a modified Hamilton function
of the kind considered long ago by
Duane et al.~\cite{DuaneEtAl}.
